%!TEX root = ../chapter07.tex 
%----------------------------------------------------------------------------------------
%	
%----------------------------------------------------------------------------------------
%----------------------------------------------------------------------------------------
%	 
%----------------------------------------------------------------------------------------


\section{Vocabulaire}
Une \textbf{équation à une inconnue} est une égalité dans laquelle apparaît une ou plusieurs lettres.

\textbf{Une solution} de l’équation est une valeur de la ou les inconnues pour lesquelles l’égalité est vraie. 
\begin{example}\hfill \\
	Soit l'équation $2x+3=x-5$ d'inconnue $x$.
	\begin{enumerate}[label=\alph*),leftmargin=*]
		\item $x=0$ n'est pas solution de l'équation car l'égalité $2\times 0+3=0-5$ est fausse
		\item $x=-8$ est une solution de l'équation, car $2\times(-8)+3=(-8)-5$ est vraie.
	\end{enumerate}
\end{example}
\begin{example}\hfill \\
	Soit l'équation $2x+3y=15$ d'inconnues $x$ et $y$.
	\begin{enumerate}[label=\alph*),leftmargin=*]
		\item Le couple  $x=6$ et $y=1$  est un couple solution car $2\times 6+3\times 1=15$.
		\item Le couple  $x=-3$ et $y=7$  est un couple solution car $2\times (-3)+3\times 7=15$.
		\item Le couple $x=1$ et $y=6$ n'est pas un couple solution car $2\times 1+3\times 6\neq 15$.  
	\end{enumerate}
\end{example}
\begin{definition} Résoudre une équation dans $\mathbb{R}$ c'est trouver toutes les valeurs  réelles des  inconnues qui rendent l'égalité vraie.
	
	On parle de résolution dans   $\mathbb{N}$ si on cherche uniquement les valeurs entières.
\end{definition} 
\begin{example}\hfill
	\begin{enumerate}[label=\alph*),leftmargin=*]
		\item L'équation $3x=1$ inconnue $x$ n'admet pas de solutions entières.
		\item L'équation $x^2+1=0$ inconnue $x$, n'admet pas de solutions dans $\mathbb{R}$.
		\item L'équation $x^2-2y^2=0$ d'inconnues $x$ et $y$ n'admet pas de solutions entières.
		\item $x=3$ et $y=2$ est un couple d'entiers solution de l'équation $x^2-2y^2=1$ d'inconnues $x$ et $y$.
	\end{enumerate}
\end{example}

\begin{definition}
	Deux équations sont dites \textbf{équivalentes} (symbole $\iff$) si elles ont le même ensemble de solutions c.à.d elles sont vraies pour les mêmes valeurs de $x$.
\end{definition} 
\begin{example}\hfill 
	\begin{enumerate}[label=\alph*),leftmargin=*]
		\item L'équations $x^2=x$ d'inconnue $x$ a pour solutions $x=0$ et $1$.\\
		L'équation $2x=x+1$ d'inconnue $x$ a une solution unique $x=1$.\\
		Les équations ne sont pas équivalentes.
		\item Les équations $2x=x+1$ et $4x=x+3$ d'inconnues $x$ ont pour seule solution  $x=1$. Elles sont équivalentes.
	\end{enumerate}
\end{example}

%\begin{theorem}\label{thm:equaequivalentes} \hfill 
%	
%	$A$, $B$ et $C$ désignent des expressions qui peuvent dépendre d'une inconnue $x$. 
%	\begin{enumerate}
%		\item  Les équations $A=B$ et $A+C=B+C$ sont équivalentes. 
%		\item  À condition que $C\neq 0$, les équations $A=B$ et $A \times C =B  \times C$ sont équivalentes.  
%	\end{enumerate} 
%\end{theorem}
%\begin{proof}\hfill 
%	
%	\vfill 
%\end{proof}

%\begin{exercise}[erreur à éviter]  \hfill 
%	\begin{DispWithArrows*}
%		(x-3) (x-1)  & =5(x-1)  \Arrow{$\divisionsymbol (x-1)$}\\
%		(x-3) & = 5  \Arrow{$+ 3$}\\
%		x & =8  
%	\end{DispWithArrows*} 
%	Les équations $(x-3) (x-1)  =5(x-1)$ et $x=8$ sont équivalentes. L'ensemble des solutions est $\mathcal{S}=\{8\}$.
%\end{exercise}